\chapter{Introduction}\label{ch:intro}

Electric vehicles depend on lithium-ion packs that are safe only inside a narrow window of voltage, current and temperature \citep{feng2018}. The battery management system (BMS) keeps the pack inside this window: it measures cell voltages, current and temperature, estimates hidden quantities such as the state of charge (SOC) and the core temperature, and decides how much current the vehicle may use.

Recent work combines the BMS with a \emph{digital twin}, a virtual replica of the battery that is kept in sync with the real one, and with cloud services for remote monitoring \citep{li2020,nejakar2025}. Nejakar et al.\ built such a system with an ESP32, low-cost sensors, a Firebase database and a web dashboard, and estimated the SOC with a hybrid Kalman filter and recurrent neural network \citep{nejakar2025}. Their system monitors and raises threshold alerts, but it does not control the battery, and it is tested on a single 12\,V battery.

This thesis takes the same three building blocks (digital twin, SOC estimation, cloud monitoring) and adds the part a vehicle needs: a BMS controller that limits the current and reacts safely to faults. The digital twin is used in a second role as well: it replaces the physical pack on the test bench, so that faults which are too dangerous for real cells can be tested repeatably.

\section*{Contributions}
\begin{enumerate}
  \item A BMS controller for a 19S12P pack on an STM32 microcontroller: SOC estimation with an extended Kalman filter (EKF), core-temperature estimation, detection of unexplained heating and a graded current limit (NORMAL, DERATE, LIMP-HOME, SAFE-STOP).
  \item A digital twin of the pack on a separate microcontroller, fitted to published measurements of the M50T cell \citep{kirkaldy2024} and connected only through CAN.
  \item Cloud monitoring through an HTTPS gateway and a web dashboard, with measured link quality.
  \item Verification from simulation to the hardware bench, including 1000 randomised runs on the microcontroller.
\end{enumerate}

\Cref{ch:background} gives the background, \cref{ch:method} the methodology, \cref{ch:hardware} the hardware, \cref{ch:results} the results and \cref{ch:conclusion} the conclusion.
